\unitopen{18}
\begin{candos}
\item podać powód i uzasadnić decyzję zdaniem z \de{weil};
\item wyrazić opinię, nadzieję i wiedzę zdaniem z \de{dass};
\item zbudować zdanie warunkowe z \de{wenn} i pytanie pośrednie z \de{ob};
\item zachować właściwy szyk: odmieniony czasownik na końcu zdania podrzędnego.
\end{candos}

\section{Słownictwo}
\vocabtable{18}

\section{Zdanie podrzędne: czasownik na końcu}
Zdanie główne ma czasownik na 2.\ miejscu; w \emph{zdaniu podrzędnym} (po \de{weil, dass, wenn, ob}) odmieniony czasownik
idzie \emph{na koniec}. Po przecinku. Przy czasowniku z klamrą (modalny, Perfekt) odmieniony element stoi \emph{ostatni}.

\begin{center}\small
\begin{tabular}{@{}ll@{}}
\toprule
\textsf{\textbf{Spójnik}} & \textsf{\textbf{Przykład}}\\
\midrule
\de{weil} (ponieważ) & \de{Ich komme später, weil der Zug Verspätung \textbf{hat}.}\\
\de{dass} (że) & \de{Ich glaube, dass er morgen \textbf{kommt}.}\\
\de{wenn} (jeśli, gdy) & \de{Wenn du Zeit \textbf{hast}, rufe ich dich an.}\\
\de{ob} (czy, pośrednie) & \de{Ich weiß nicht, ob er schon \textbf{angekommen ist}.}\\
\de{denn} (bo) & \de{Ich komme später, denn der Zug \textbf{hat} Verspätung.} (szyk jak w głównym)\\
\bottomrule
\end{tabular}
\end{center}
Gdy zdanie podrzędne \emph{poprzedza} główne, czasownik głównego idzie na pierwsze miejsce po przecinku:
\de{Wenn du Zeit hast, \textbf{rufe} ich dich an.} (odmieniony \de{rufe} przed podmiotem \de{ich}).

\begin{zasada}[Trzy testy]
\textbf{(1)} Czy po spójniku czasownik jest na końcu? \textbf{(2)} Czy zdanie podrzędne jest oddzielone przecinkiem?
\textbf{(3)} Jeśli zaczynasz od podrzędnego, czy następny wyraz to czasownik? Przy \de{ob}: pytanie tak/nie zamienia
się na zdanie podrzędne (\de{Kommt er? $\rightarrow$ Ich weiß nicht, ob er kommt.}).
\end{zasada}

\begin{falle}
Polskie „bo” i „ponieważ” nie zmieniają szyku, więc nawyk zostaje: \de{Ich bleibe zu Hause, weil ich bin krank}
(błąd). Poprawnie: \de{\dots, weil ich krank \textbf{bin}.} Drugi: \de{weil} w mowie potocznej bywa łączone z
szykiem głównym (\de{weil ich bin krank}), ale w egzaminie i w piśmie jest to błąd. Trzeci: nie myl \de{dass}
(spójnik) z \de{das} (zaimek lub rodzajnik): \de{Ich weiß, \textbf{dass} das Auto kaputt ist.}
\end{falle}

\begin{szwajcaria}
Szwajcarzy w mowie omijają zdania podrzędne i mówią krótko; w e-mailach używają ich jednak regularnie.
Uprzejme uzasadnienie to \de{Leider kann ich nicht kommen, weil ich einen anderen Termin habe.} W Szwajcarii
nie wypada jedynie odmawiać: dodaje się powód.
\end{szwajcaria}

\begin{praca}
\de{Ich glaube, dass der Fehler im Cache liegt.} \de{Ich weiß nicht, ob der Test grün ist.} \de{Wenn der Build fehlschlägt,
schicke ich dir die Logs.} \de{Ich kann heute nicht kommen, weil ich krank bin.} Takie zdania są rdzeniem
standupów i opisów PR.
\end{praca}

\section{Dialogi}
\textbf{Dialog 1 --- spóźnienie}
\begin{dialog}
\turn{A}{Warum bist du zu spät?}{Dlaczego się spóźniłeś?}
\turn{B}{Entschuldigung. Ich bin zu spät, weil der Zug Verspätung hatte.}{Przepraszam. Spóźniłem się, bo pociąg miał opóźnienie.}
\turn{A}{Kein Problem. Wir haben noch nicht angefangen, weil der Chef fehlt.}{Nie ma problemu. Jeszcze nie zaczęliśmy, bo brakuje szefa.}
\turn{B}{Weißt du, ob er heute kommt?}{Czy wiesz, czy on dziś przyjdzie?}
\turn{A}{Ich glaube, dass er im Homeoffice arbeitet.}{Myślę, że pracuje zdalnie.}
\end{dialog}

\section{Ćwiczenia}
\begin{cwiczenia}
\exitem{Połącz: \de{Ich bleibe zu Hause. Ich bin krank.} (weil)}{\de{Ich bleibe zu Hause, weil ich krank bin.}}
\exitem{Połącz: \de{Ich glaube. Er kommt morgen.} (dass)}{\de{Ich glaube, dass er morgen kommt.}}
\exitem{Połącz: \de{Du hast Zeit. Ich rufe dich an.} (wenn, zaczynając od wenn)}{\de{Wenn du Zeit hast, rufe ich dich an.}}
\exitem{Zamień na pytanie pośrednie: \de{Kommt er?} $\rightarrow$ \de{Ich weiß nicht, \dots}}{\de{Ich weiß nicht, ob er kommt.}}
\exitem{Zamień na pytanie pośrednie: \de{Ist der Test grün?} $\rightarrow$ \de{Ich weiß nicht, \dots}}{\de{Ich weiß nicht, ob der Test grün ist.}}
\exitem{Popraw: \de{Ich bleibe hier, weil ich bin müde.}}{\de{Ich bleibe hier, weil ich müde bin.}}
\exitem{Popraw: \de{Ich weiß, das er kommt.}}{\de{Ich weiß, dass er kommt.}}
\exitem{Perfekt w zdaniu podrzędnym: \de{Ich weiß, dass er schon \dots{}} (ankommen)}{\de{Ich weiß, dass er schon angekommen ist.}}
\exitem{Modalne w zdaniu podrzędnym: \de{Er kommt nicht, weil er arbeiten \dots{}} (müssen)}{\de{Er kommt nicht, weil er arbeiten muss.}}
\exitem{Użyj \de{denn} zamiast \de{weil}: \de{Ich bleibe zu Hause, weil ich krank bin.}}{\de{Ich bleibe zu Hause, denn ich bin krank.}}
\exitem{Przetłumacz: „Nie wiem, czy to działa.”}{\de{Ich weiß nicht, ob das funktioniert.}}
\exitem{Przetłumacz: „Spóźnię się, bo mam spotkanie.”}{\de{Ich komme später, weil ich ein Meeting habe.}}
\end{cwiczenia}

\section{Trening słownictwa (generowany)}
\drill{18}

\weekend{Napisz osiem zdań o swoim dniu, z użyciem po dwa razy \de{weil}, \de{dass}, \de{wenn}, \de{ob}.
Zaznacz kolorem czasowniki na końcu. Wyślij do korekty szyku.}
